\section{BUGHUNTER}


Utilizamos el Conjunto de datos BugHunter: un tipo novedoso de conjunto de datos de errores construido automáticamente y disponible gratuitamente que contiene elementos de código (archivos, clases, métodos) con un amplio conjunto de métricas de código e información de errores \cite{ferenc_automatically_2020}.

Otros conjuntos de datos de errores disponibles siguen el enfoque tradicional (PROMISE \cite{sayyad_shirabad_promise_2005}, Eclipse, entre otros) de recopilar las características de todos los elementos del código fuente (con errores y sin errores) en solo una o más versiones de lanzamiento preseleccionadas del código. El enfoque seleccionado, por otro lado, captura el error y los estados fijos de los mismos elementos del código fuente desde el marco de tiempo más estrecho que podemos identificar para detectar la presencia de un error, independientemente de las versiones de lanzamiento.

El conjunto de datos contiene información útil de caracterización de elementos de código defectuoso, se utilizan varios proyectos almacenados en GitHub que son mantenidos regularmente. Se utiliza la API de GitHub que hace posible implementar una herramienta de recuperación automática de datos. Se seleccionan 15 proyectos Java como sistemas temáticos, que difieren en muchos aspectos entre sí (tamaño, dominio, número de errores informados) para cubrir un conjunto amplio y general de sistemas \cite{ferenc_automatically_2020}. 

El conjunto de datos se basa en un nuevo enfoque que recopila información del instante antes y después de la corrección de los elementos del código fuente (y sus características) que se vieron afectados por errores y no considera aquellos elementos del código que no fueron afectados por errores. Este tipo de conjunto de datos ayuda a capturar los cambios en las métricas de los productos de software cuando se corrige un error. Por lo tanto, se puede aprender de las diferencias en las métricas del código fuente entre elementos de código defectuosos y no defectuosos.

En la tabla \ref{table:datasetmeta}, resumimos los metadatos principales del conjunto de datos. La cantidad de métricas existentes en cada categoría y la cantidad de registros existentes en cada uno de los niveles de granularidad.

\begin{table}[H]
\captionsetup{labelsep=period, labelfont=bf}
\caption{Metadatos del conjunto de datos}
\centering
\begin{tabular}{lc}
\toprule
\textbf{Granularity}        & file, class, method                                                                                                                \\\midrule
\textbf{Number of projects} & 15                                                                                                                                 \\ 
\hline
\textbf{Number of metrics}  & \begin{tabular}[c]{@{}c@{}}static source code metrics (66),\\ code duplication metrics (8),\\ code smell metrics (35)\end{tabular} \\
\hline
\textbf{Number of entries}  & \begin{tabular}[c]{@{}c@{}}file: 70,088\\ class: 84,562\\ method: 159,078


\end{tabular}  
\\\bottomrule
\end{tabular}

\label{table:datasetmeta}
\end{table} 

 Realizamos una descripción de los datos en \ref{description-table} mostrando cantidad, media, desviación estándar, mínimo, máximo y cantidad de elementos. Al analizar la tabla descriptiva de los datos podemos concluir:
 {
\fontsize{14pt}{14pt}\selectfont
\begin{itemize}
\item No existen datos nulos dado que el contados de todos es igual al total de registros.
\item Gran variedad en los rangos de los datos.
\item Una gran cantidad de características que tienen moda y mediana iguales y con valor 0, pudiendo implicar poca variación en sus datos.
\end{itemize}
}

El conjunto de datos BugHunter cuenta con varias métricas de código extraídas con la herramienta OpenStaticAnalyzer \cite{openstaticanalyzergithubio} utilizando el código etiquetado en el GitHub de 15 proyectos, escritos en Java. Especialmente los más grandes y que tengan una cantidad adecuada de problemas etiquetados como errores, porque son más adecuados para este tipo de análisis.

\begin{table}[H]
{
\captionsetup{labelsep=period, labelfont=bf}
\caption{ Description Table}
\fontsize{9pt}{6pt}\selectfont
\begin{tabular}{@{}lp{10mm}p{10mm}p{10mm}p{10mm}llllll@{}} 
\toprule
Medida& Media & Error   típico & Mediana & Moda & Desviación   estándar & Varianza   de la muestra & Mínimo & Máximo & Suma& Cuenta \\ 
\midrule
CC                                  & 0,12           & 0,00                    & 0,00             & 0,00          & 0,30                           & 0,09                              & 0,00            & 1,00            & 19334,08       & 159078,00       \\
CCL                                 & 0,66           & 0,03                    & 0,00             & 0,00          & 11,53                          & 133,02                            & 0,00            & 499,00          & 104713,00      & 159078,00       \\
CCO                                 & 4,01           & 0,71                    & 0,00             & 0,00          & 282,97                         & 80069,76                          & 0,00            & 75248,00        & 638669,00      & 159078,00       \\
CI                                  & 1,95           & 0,14                    & 0,00             & 0,00          & 56,22                          & 3160,84                           & 0,00            & 4671,00         & 310335,00      & 159078,00       \\
CLC                                 & 0,12           & 0,00                    & 0,00             & 0,00          & 0,30                           & 0,09                              & 0,00            & 1,00            & 19071,79       & 159078,00       \\
LDC                                 & 7,52           & 0,37                    & 0,00             & 0,00          & 146,09                         & 21343,22                          & 0,00            & 6691,00         & 1196598,00     & 159078,00       \\
LLDC                                & 6,72           & 0,33                    & 0,00             & 0,00          & 132,41                         & 17531,44                          & 0,00            & 5848,00         & 1068422,00     & 159078,00       \\
HCPL                                & 202,83         & 1,94                    & 106,09           & 13,61         & 771,92                         & 595862,35                         & 2,00            & 57839,10        & 32266128,82    & 159078,00       \\
HDIF                                & 26,30          & 0,96                    & 15,04            & 3,50          & 384,17                         & 147587,63                         & 1,00            & 56556,50        & 4183142,22     & 159078,00       \\
HEFF                                & 250671,02      & 65265,39                & 3554,80          & 49,13         & 26030828,33                    & 677604023793090,00                & 4,75            & 3944690000,00   & 39876245063,67 & 159078,00       \\
HNDB                                & 910,90         & 43,12                   & 232,92           & 13,41         & 17198,86                       & 295800753,28                      & 2,83            & 2496560,00      & 144904764,53   & 159078,00       \\
HPL                                 & 120,57         & 2,12                    & 48,00            & 12,00         & 844,70                         & 713519,30                         & 3,00            & 39378,00        & 19180331,00    & 159078,00       \\
HPV                                 & 39,73          & 0,19                    & 27,00            & 10,00         & 74,43                          & 5539,97                           & 3,00            & 4764,00         & 6319883,00     & 159078,00       \\
HTRP                                & 13926,17       & 3625,86                 & 197,49           & 2,73          & 1446157,65                     & 2091371943879,64                  & 0,26            & 219150000,00    & 2215347324,87  & 159078,00       \\
HVOL                                & 848,92         & 23,68                   & 230,00           & 19,65         & 9444,64                        & 89201173,08                       & 4,75            & 451128,00       & 135044408,26   & 159078,00       \\
MI                                  & 104,36         & 0,06                    & 106,10           & 132,83        & 23,81                          & 567,01                            & -579,07         & 162,66          & 16602174,03    & 159078,00       \\
MIMS                                & 61,05          & 0,03                    & 62,05            & 77,68         & 13,74                          & 188,78                            & 0,00            & 95,12           & 9712500,44     & 159078,00       \\
MISEI                               & 83,11          & 0,08                    & 84,46            & 116,03        & 33,76                          & 1139,90                           & -668,08         & 208,70          & 13221672,48    & 159078,00       \\
MISM                                & 48,73          & 0,05                    & 49,39            & 67,85         & 19,25                          & 370,44                            & 0,00            & 122,05          & 7751585,71     & 159078,00       \\
McCC                                & 3,54           & 0,03                    & 1,00             & 1,00          & 12,08                          & 146,02                            & 1,00            & 2387,00         & 562498,00      & 159078,00       \\
NL                                  & 1,08           & 0,01                    & 0,00             & 0,00          & 2,03                           & 4,12                              & 0,00            & 82,00           & 171181,00      & 159078,00       \\
NLE                                 & 0,94           & 0,00                    & 0,00             & 0,00          & 1,35                           & 1,82                              & 0,00            & 27,00           & 150219,00      & 159078,00       \\
NII                                 & 1,60           & 0,03                    & 0,00             & 0,00          & 13,08                          & 171,14                            & 0,00            & 1833,00         & 253982,00      & 159078,00       \\
NOI                                 & 5,79           & 0,02                    & 3,00             & 1,00          & 8,50                           & 72,30                             & 0,00            & 185,00          & 920814,00      & 159078,00       \\
CD                                  & 0,05           & 0,00                    & 0,00             & 0,00          & 0,13                           & 0,02                              & 0,00            & 0,95            & 8743,04        & 159078,00       \\
CLOC                                & 1,35           & 0,01                    & 0,00             & 0,00          & 3,81                           & 14,53                             & 0,00            & 198,00          & 214253,00      & 159078,00       \\
DLOC                                & 0,54           & 0,01                    & 0,00             & 0,00          & 2,39                           & 5,70                              & 0,00            & 195,00          & 86694,00       & 159078,00       \\
TCD                                 & 0,05           & 0,00                    & 0,00             & 0,00          & 0,13                           & 0,02                              & 0,00            & 0,95            & 8603,51        & 159078,00       \\
TCLOC                               & 1,40           & 0,01                    & 0,00             & 0,00          & 3,92                           & 15,36                             & 0,00            & 198,00          & 222238,00      & 159078,00       \\
LLOC                                & 19,96          & 0,32                    & 9,00             & 4,00          & 127,48                         & 16250,40                          & 1,00            & 5851,00         & 3175814,00     & 159078,00       \\
LOC                                 & 22,73          & 0,35                    & 10,00            & 4,00          & 141,47                         & 20015,14                          & 1,00            & 6273,00         & 3615214,00     & 159078,00       \\
NOS                                 & 13,33          & 0,27                    & 5,00             & 1,00          & 106,34                         & 11308,68                          & 0,00            & 5079,00         & 2120731,00     & 159078,00       \\
NUMPAR                              & 1,14           & 0,00                    & 1,00             & 0,00          & 1,60                           & 2,56                              & 0,00            & 38,00           & 180569,00      & 159078,00       \\
TLLOC                               & 21,51          & 0,34                    & 10,00            & 4,00          & 134,48                         & 18085,43                          & 1,00            & 5851,00         & 3421101,00     & 159078,00       \\
TLOC                                & 24,43          & 0,37                    & 11,00            & 4,00          & 148,75                         & 22125,86                          & 1,00            & 6703,00         & 3885847,00     & 159078,00       \\
TNOS                                & 14,01          & 0,27                    & 5,00             & 1,00          & 108,60                         & 11794,42                          & 0,00            & 5079,00         & 2229242,00     & 159078,00       \\
WarningBlocker                      & 0,00           & 0,00                    & 0,00             & 0,00          & 0,00                           & 0,00                              & 0,00            & 1,00            & 1,00           & 159078,00       \\
WarningCritical                     & 0,11           & 0,00                    & 0,00             & 0,00          & 0,81                           & 0,66                              & 0,00            & 45,00           & 17572,00       & 159078,00       \\
WarningInfo                         & 0,00           & 0,00                    & 0,00             & 0,00          & 0,00                           & 0,00                              & 0,00            & 0,00            & 0,00           & 159078,00       \\
WarningMajor                        & 0,49           & 0,00                    & 0,00             & 0,00          & 1,80                           & 3,26                              & 0,00            & 79,00           & 78578,00       & 159078,00       \\
WarningMinor                        & 2,32           & 0,10                    & 0,00             & 0,00          & 39,21                          & 1537,67                           & 0,00            & 2889,00         & 368755,00      & 159078,00       \\
Android Rules                       & 0,00           & 0,00                    & 0,00             & 0,00          & 0,00                           & 0,00                              & 0,00            & 0,00            & 0,00           & 159078,00       \\
Basic Rules                         & 0,05           & 0,00                    & 0,00             & 0,00          & 0,42                           & 0,18                              & 0,00            & 34,00           & 7199,00        & 159078,00       \\
Brace Rules                         & 0,35           & 0,02                    & 0,00             & 0,00          & 9,89                           & 97,87                             & 0,00            & 2290,00         & 55578,00       & 159078,00       \\
Clone Implementation Rules          & 0,00           & 0,00                    & 0,00             & 0,00          & 0,02                           & 0,00                              & 0,00            & 2,00            & 55,00          & 159078,00       \\
Code Size Rules                     & 0,00           & 0,00                    & 0,00             & 0,00          & 0,00                           & 0,00                              & 0,00            & 0,00            & 0,00           & 159078,00       \\
Comment Rules                       & 0,00           & 0,00                    & 0,00             & 0,00          & 0,00                           & 0,00                              & 0,00            & 0,00            & 0,00           & 159078,00       \\
Controversial Rules                 & 0,06           & 0,00                    & 0,00             & 0,00          & 0,26                           & 0,07                              & 0,00            & 11,00           & 10096,00       & 159078,00       \\
Coupling Rules                      & 0,00           & 0,00                    & 0,00             & 0,00          & 0,00                           & 0,00                              & 0,00            & 0,00            & 0,00           & 159078,00       \\
Design Rules                        & 1,41           & 0,09                    & 0,00             & 0,00          & 36,77                          & 1351,85                           & 0,00            & 2807,00         & 224629,00      & 159078,00       \\
Empty Code Rules                    & 0,02           & 0,00                    & 0,00             & 0,00          & 0,25                           & 0,06                              & 0,00            & 34,00           & 2414,00        & 159078,00       \\
Finalizer Rules                     & 0,00           & 0,00                    & 0,00             & 0,00          & 0,01                           & 0,00                              & 0,00            & 1,00            & 24,00          & 159078,00       \\
Import Statement Rules              & 0,06           & 0,00                    & 0,00             & 0,00          & 1,48                           & 2,19                              & 0,00            & 81,00           & 9044,00        & 159078,00       \\
J2EE Rules                          & 0,00           & 0,00                    & 0,00             & 0,00          & 0,04                           & 0,00                              & 0,00            & 5,00            & 132,00         & 159078,00       \\
JUnit Rules                         & 0,43           & 0,01                    & 0,00             & 0,00          & 2,02                           & 4,07                              & 0,00            & 150,00          & 68963,00       & 159078,00       \\
Jakarta Commons Logging Rules       & 0,01           & 0,00                    & 0,00             & 0,00          & 0,18                           & 0,03                              & 0,00            & 9,00            & 2056,00        & 159078,00       \\
Java Logging Rules                  & 0,04           & 0,00                    & 0,00             & 0,00          & 0,39                           & 0,15                              & 0,00            & 46,00           & 6637,00        & 159078,00       \\
JavaBean Rules                      & 0,00           & 0,00                    & 0,00             & 0,00          & 0,00                           & 0,00                              & 0,00            & 0,00            & 0,00           & 159078,00       \\
MigratingToJUnit4 Rules             & 0,00           & 0,00                    & 0,00             & 0,00          & 0,00                           & 0,00                              & 0,00            & 0,00            & 0,00           & 159078,00       \\
Migration Rules                     & 0,00           & 0,00                    & 0,00             & 0,00          & 0,04                           & 0,00                              & 0,00            & 6,00            & 124,00         & 159078,00       \\
Migration13 Rules                   & 0,00           & 0,00                    & 0,00             & 0,00          & 0,00                           & 0,00                              & 0,00            & 0,00            & 0,00           & 159078,00       \\
Migration14 Rules                   & 0,00           & 0,00                    & 0,00             & 0,00          & 0,00                           & 0,00                              & 0,00            & 0,00            & 0,00           & 159078,00       \\
Migration15 Rules                   & 0,00           & 0,00                    & 0,00             & 0,00          & 0,00                           & 0,00                              & 0,00            & 0,00            & 0,00           & 159078,00       \\
Naming Rules                        & 0,05           & 0,00                    & 0,00             & 0,00          & 0,25                           & 0,06                              & 0,00            & 8,00            & 8554,00        & 159078,00       \\
Optimization Rules                  & 0,01           & 0,00                    & 0,00             & 0,00          & 0,16                           & 0,02                              & 0,00            & 15,00           & 1787,00        & 159078,00       \\
Security Code Guideline Rules       & 0,00           & 0,00                    & 0,00             & 0,00          & 0,06                           & 0,00                              & 0,00            & 3,00            & 552,00         & 159078,00       \\
Strict Exception Rules              & 0,11           & 0,00                    & 0,00             & 0,00          & 0,44                           & 0,19                              & 0,00            & 14,00           & 17806,00       & 159078,00       \\
String and StringBuffer Rules       & 0,11           & 0,00                    & 0,00             & 0,00          & 0,97                           & 0,94                              & 0,00            & 88,00           & 17568,00       & 159078,00       \\
Type Resolution Rules               & 0,16           & 0,00                    & 0,00             & 0,00          & 0,41                           & 0,17                              & 0,00            & 11,00           & 24819,00       & 159078,00       \\
Unnecessary and   Unused Code Rules & 0,06           & 0,00                    & 0,00             & 0,00          & 1,45                           & 2,11                              & 0,00            & 79,00           & 9382,00        & 159078,00       \\
Vulnerability Rules                 & 0,00           & 0,00                    & 0,00             & 0,00          & 0,00                           & 0,00                              & 0,00            & 1,00            & 1,00           & 159078,00       \\
Number of Bugs                      & 0,47           & 0,00                    & 0,00             & 0,00          & 0,70                           & 0,50                              & 0,00            & 6,00            & 75096,00       & 159078,00       \\ 
\bottomrule
\end{tabular}
\label{description-table}
}
\end{table}











En la \ref{metrics-table} se muestra al conjunto de métricas y sus niveles correspondientes, el BugHunter contiene las métricas de código en tres niveles (método, clase y archivo).

\begin{table}[H]
{
\fontsize{8pt}{4pt}\selectfont
\captionsetup{labelsep=period, labelfont=bf}
\caption{Métricas del código fuente}
\centering
\begin{tabular}{lllll}
\toprule
Abreviación & Nombre Completo                               & Método & Clase & Archivo \\
\midrule
CLOC        & Comment Lines of Code                         & X      & X     & X       \\
LOC         & Lines of Code                                 & X      & X     & X       \\
LLOC        & Logical Lines of Code                         & X      & X     & X       \\
NL          & Nesting Level                                 & X      & X     &         \\
NLE         & Nesting Level Else-If                         & X      & X     &         \\
NII         & Number of Incoming   Invocations              & X      & X     &         \\
NOI         & Number of Outgoing Invocations                & X      & X     &         \\
CD          & Comment Density                               & X      & X     &         \\
DLOC        & Documentation Lines of Code                   & X      & X     &         \\
TCD         & Total Comment Density                         & X      & X     &         \\
TCLOC       & Total Comment   Lines of Code                 & X      & X     &         \\
NOS         & Number of Statements                          & X      & X     &         \\
TLOC        & Total Lines of Code                           & X      & X     &         \\
TLLOC       & Total   Logical Lines of Code                 & X      & X     &         \\
TNOS        & Total Number of Statements                    & X      & X     &         \\
McCC        & McCabe’s Cyclomatic Comple   ity              & X      &       & X       \\
PDA         & Public Documented API                         &        & X     & X       \\
PUA         & Public Undocumented API                       &        & X     & X       \\
HCPL        & Halstead Calculated Program Length            & X      &       &         \\
HDIF        & Halstead Difficulty                           & X      &       &         \\
HEFF        & Halstead Effort                               & X      &       &         \\
HNDB        & Halstead   Number of Delivered Bugs           & X      &       &         \\
HPL         & Halstead Program Length                       & X      &       &         \\
HPV         & Halstead Program   Vocabulary                 & X      &       &         \\
HTRP        & Halstead Time   Required to Program           & X      &       &         \\
HVOL        & Halstead Volume                               & X      &       &         \\
MIMS        & Maintainability Index (Microsoft version)     & X      &       &         \\
MI          & Maintainability Index   (Original version)    & X      &       &         \\
MISEI       & Maintainability Index (SEI version)           & X      &       &         \\
MISM        & Maintainability Index   (SourceMeter version) & X      &       &         \\
NUMPAR      & Number of Parameters                          & X      &       &         \\
LCOM5       & Lack   of Cohesion in Methods                 &        & X     &         \\
WMC         & Weighted Methods per Class                    &        & X     &         \\
CBO         & Coupling Between Object   classes             &        & X     &         \\
CBOI        & Coupling Between   Object classes Inverse     &        & X     &         \\
RFC         & Response set For Class                        &        & X     &         \\
AD          & API Documentation                             &        & X     &         \\
DIT         & Depth of Inheritance   Tree                   &        & X     &         \\
NOA         & Number of Ancestors                           &        & X     &         \\
NOC         & Number of Children                            &        & X     &         \\
NOD         & Number of Descendants                         &        & X     &         \\
NOP         & Number of Parents                             &        & X     &         \\
NA          & Number of Attributes                          &        & X     &         \\
NG          & Number of Getters                             &        & X     &         \\
NLA         & Number of Local Attributes                    &        & X     &         \\
NLG         & Number of Local Getters                       &        & X     &         \\
NLM         & Number of Local Methods                       &        & X     &         \\
NLPA        & Number   of Local Public Attributes           &        & X     &         \\
NLPM        & Number of Local   Public Methods              &        & X     &         \\
NLS         & Number of Local Setters                       &        & X     &         \\
NM          & Number of Methods                             &        & X     &         \\
NPA         & Number of Public   Attributes                 &        & X     &         \\
NPM         & Number of Public Methods                      &        & X     &         \\
NS          & Number of Setters                             &        & X     &         \\
TNA         & Total Number of Attributes                    &        & X     &         \\
TNG         & Total Number of Getters                       &        & X     &         \\
TNLA        & Total Number of   Local Attributes            &        & X     &         \\
TNLG        & Total   Number of Local Getters               &        & X     &         \\
TNLM        & Total Number of   Local Methods               &        & X     &         \\
TNLPA       & Total   Number of Local Public Attributes     &        & X     &         \\
TNLPM       & Total Number of   Local Public Methods        &        & X     &         \\
TNLS        & Total   Number of Local Setters               &        & X     &         \\
TNM         & Total Number of Methods                       &        & X     &         \\
TNPA        & Total   Number of Public Attributes           &        & X     &         \\
TNPM        & Total Number of   Public Methods              &        & X     &         \\
TNS         & Total Number of Setters                       &        & X     &        \\
\bottomrule
\end{tabular}

\label{metrics-table}
}
\end{table}


OpenStaticAnalyzer también proporciona un módulo de detección de violaciones de reglas de codificación mostradas en la tabla \ref{table:smell_metrics}. La presencia de violaciones de reglas en un elemento de código fuente puede causar errores en una fase posterior, por lo tanto, la cantidad de violaciones de reglas diferentes ubicadas en un elemento del código fuente puede servir como predictores y el conjunto de datos también encapsula esta información. 


\begin{table}[H]
{
\fontsize{9pt}{5pt}\selectfont
\captionsetup{labelsep=period, labelfont=bf}
\caption{Métricas de violación de reglas}
\centering
\begin{tabular}{ll}
 \toprule
\textbf{Abreviación} & \textbf{Nombre completo}          \\
\midrule
WI                   & WarningInfo                                           \\
WMA                  & WarningMajor                                         \\
AR                   & Android Rules                                          \\
BAR                  & Basic Rules                                            \\
CSR                  & Controversial Rules                                    \\
CR                   & Coupling Rules                                        \\
NR                   & Naming Rules                                           \\
OR                   & Optimization Rules                                     \\
SCGR                 & Security Code Guideline Rules                          \\
SER                  & Strict Exception Rules                                 \\
SSR                  & String and StringBuffer Rules                           \\
TRR                  & Type Resolution Rules                                  \\
VR                   & Vulnerability Rules                                   \\
J2EER                & J2EE Rules                                              \\
WC                   & WarningCritical                                          \\
UUCR                 & Unnecessary and Unused Code Rules                         \\
JLR                  & Java Logging Rules                                       \\
ISR                  & Import Statement Rules                                    \\
M15R                 & Migration15 Rules                                         \\
WMI                  & WarningMinor                                              \\
MR                   & Migration Rules                                           \\
M13R                 & Migration13 Rules                                         \\
M14R                 & Migration14 Rules                                        \\
MUR                  & MigratingToJUnit4 Rules                                   \\
JBR                  & JavaBean Rules                                            \\
JCLR                 & Jakarta Commons Logging Rules                             \\
JUR                  & JUnit Rules                                              \\
FR                   & Finalizer Rules                                           \\
ECR                  & Empty Code Rules                                          \\
DR                   & Design Rules                                              \\
CR                   & Comment Rules                                             \\
CIR                  & Clone Implementation Rules                               \\
BRR                  & Brace Rules   \\
\bottomrule
\end{tabular}

\label{table:smell_metrics}
}
\end{table}